%%%% source: https://espanol.libretexts.org/Bookshelves/Quimica/Qu%C3%ADmica_Introductoria%2C_Conceptual_y_GOB/Qu%C3%ADmica_Introductoria_(CK-12)/17%3A_Termoqu%C3%ADmica/17.04%3A_Capacidad_calor%C3%ADfica_y_calor_espec%C3%ADfico#:~:text=Resumen%20*%20La%20capacidad%20calor%C3%ADfica%20es%20la,la%20sustancia%20por%201%20o%20C%20.
%%%% Tabla de calores específicos
\begin{table}[h!]
      \centering
      \caption{calores específicos de algunas sustancias comunes}
      \label{tab:calores_especificos}
      \begin{tabular}{l c}
          \hline
          \textbf{sustancia} & \textbf{calor específico} $\left(\mathrm{j\,g^{-1}\,^{\circ}c^{-1}}\right)$ \\
          \hline
          agua (l)                & 4.18  \\
          agua (s)                & 2.06  \\
          agua (g)                & 1.87  \\
          amoníaco (g)            & 2.09  \\
          etanol (l)              & 2.44  \\
          aluminio (s)            & 0.897 \\
          carbono, grafito (s)    & 0.709 \\
          cobre (s)               & 0.385 \\
          oro (s)                 & 0.129 \\
          hierro (s)              & 0.449 \\
          plomo (s)               & 0.129 \\
          mercurio (l)            & 0.140 \\
          plata (s)               & 0.233 \\
          \hline
          \multicolumn{2}{c}{\textbf{gases a temperatura ambiente}} \\
          \hline
          nitrógeno (g)           & 1.04  \\
          oxígeno (g)             & 0.918 \\
          dióxido de carbono (g)  & 0.844 \\
          hidrógeno (g)           & 14.3  \\
          helio (g)               & 5.19  \\
          neón (g)                & 1.03  \\
          argón (g)               & 0.520 \\
          metano (g)              & 2.22  \\
          \hline
      \end{tabular}
\end{table}
